\chapter{Scientific Biology}\label{ch:g10:scientific-biology}

\omterm{def:g10:scientific-biology:science}{Biology} is not a catalogue of facts about living things: it is a way
of asking questions and checking answers against the living world.
This chapter installs the habits of scientific inquiry --- \omterm{def:g10:scientific-biology:observation}{observation},
\omterm{def:g10:scientific-biology:hypothesis}{hypothesis}, \omterm{def:g10:scientific-biology:experiment}{experiment}, control --- and names the shared properties that
mark something as alive. Every later chapter rests on these tools.

\section{What science does}

\begin{definition}[Science and biology]\label{def:g10:scientific-biology:science}
\emph{Science}\index{science} is a public method for building reliable
knowledge about the natural world: claims are stated so they can be
tested, and they stand only as long as evidence supports them.
\emph{Biology}\index{biology} is the science of life --- of living
\omterm{def:g10:scientific-biology:organism}{organisms}, their structures, functions, interactions and history.
\end{definition}

\begin{definition}[Observation and inference]\label{def:g10:scientific-biology:observation}
An \emph{observation}\index{observation} is a record of what can be
sensed or measured, with instruments when needed: colour, mass, count,
temperature, DNA sequence. An \emph{inference}\index{inference} is a
reasoned conclusion drawn from observations; it may be wrong even when
the observations are correct.
\end{definition}

\begin{example}[Leaves and light]\label{ex:g10:scientific-biology:obs}
You observe that the underside of a forest leaf is paler than the
upper side. A possible \omterm{def:g10:scientific-biology:observation}{inference} is that less chlorophyll sits there;
another is that the cuticle is thinner. Only further measurement
(pigment extraction, microscopy) can decide.
\end{example}

\begin{definition}[Hypothesis]\label{def:g10:scientific-biology:hypothesis}
A \emph{hypothesis}\index{hypothesis} is a provisional, testable
explanation of a pattern in nature. It must make predictions that
could fail: if no \omterm{def:g10:scientific-biology:observation}{observation} could prove it wrong, it is not a
scientific hypothesis.
\end{definition}

\begin{remark}[Hypothesis versus guess]\label{rem:g10:scientific-biology:guess}
``I think plants need light'' is a vague hope. ``If light is required
for net sugar production, then seedlings kept in continuous darkness
will not gain dry mass'' is a \omterm{def:g10:scientific-biology:hypothesis}{hypothesis}: it names a comparison and
an outcome that could go either way.
\end{remark}

\section{Experiment, variables and controls}

\begin{definition}[Experiment]\label{def:g10:scientific-biology:experiment}
An \emph{experiment}\index{experiment} is a planned comparison in which
one factor is deliberately changed while others are held as steady as
possible, so that a measured response can be attributed to that factor.
\end{definition}

\begin{definition}[Variables]\label{def:g10:scientific-biology:variables}
In an \omterm{def:g10:scientific-biology:experiment}{experiment} the \emph{independent variable}\index{variable!independent}
is the factor the investigator sets (for example light intensity). The
\emph{dependent variable}\index{variable!dependent} is the response that
is measured (for example \omterm{def:g10:scientific-biology:properties}{growth} rate). \emph{Controlled
variables}\index{variable!controlled} are factors kept the same so they
do not cloud the comparison (temperature, water, soil type).
\end{definition}

\begin{omfigure}
\begin{tikzpicture}[font=\small,
  box/.style={draw, rounded corners, fill=omDef!10, minimum width=2.8cm,
    minimum height=1.0cm, align=center},
  arr/.style={-Stealth, thick, omThm}]
  \node[box] (iv) at (0,0) {independent\\variable};
  \node[box] (sys) at (4.0,0) {living\\system};
  \node[box] (dv) at (8.0,0) {dependent\\variable};
  \node[box, fill=omMeth!12] (cv) at (4.0,-2.0) {controlled\\variables};
  \draw[arr] (iv) -- (sys);
  \draw[arr] (sys) -- (dv);
  \draw[arr, omMeth] (cv) -- (sys);
\end{tikzpicture}

{\small The experimental skeleton: change one factor, measure one
response, hold the rest fixed.}
\end{omfigure}

\begin{definition}[Control group]\label{def:g10:scientific-biology:control}
A \emph{control group}\index{control group} (or control treatment)
receives every condition of the \omterm{def:g10:scientific-biology:experiment}{experiment} \emph{except} the change
under test. It is the baseline against which the experimental group
is compared. Without a control, ``it grew'' cannot mean ``it grew
\emph{because} of the treatment''.
\end{definition}

\begin{method}[Designing a fair test]\label{met:g10:scientific-biology:fair-test}
\begin{enumerate}
  \item State a clear question and a testable \omterm{def:g10:scientific-biology:hypothesis}{hypothesis}.
  \item Name the independent and \omterm{def:g10:scientific-biology:variables}{dependent variables}, with units.
  \item List controlled variables and how each will be held fixed.
  \item Include a control treatment (or a zero level of the factor).
  \item Replicate: several individuals or trials, not one lucky case.
  \item Record raw data, then summarise (mean, range) and compare.
\end{enumerate}
\end{method}

\begin{example}[Does fertiliser speed seedling growth?]\label{ex:g10:scientific-biology:fertiliser}
\omterm{def:g10:scientific-biology:hypothesis}{Hypothesis}: extra nitrogen increases mean height after \qty{14}{d}.
\omterm{def:g10:scientific-biology:variables}{Independent variable}: nitrogen dose (\unit{mg/L} in water). \omterm{def:g10:scientific-biology:variables}{Dependent
variable}: height (\unit{mm}). Controlled: light, temperature, pot size,
seed stock, water volume. Control: water with no nitrogen. Ten pots
per treatment. Measure all pots; compare means.
\end{example}

\begin{remark}[Replication is not decoration]\label{rem:g10:scientific-biology:replicate}
Living systems vary. One tall plant under fertiliser proves nothing;
ten plants on average taller than ten controls is evidence. \omterm{def:g10:scientific-biology:science}{Biology}
counts individuals for a reason.
\end{remark}

\section{From data to knowledge}

\begin{definition}[Theory and law (working sense)]\label{def:g10:scientific-biology:theory}
A scientific \emph{theory}\index{theory (scientific)} is a well-tested
framework that explains a broad range of \omterm{def:g10:scientific-biology:observation}{observations} (for example \omterm{prop:g10:scientific-biology:cell-unit}{cell
theory}, evolution by natural selection). A \emph{law}\index{law
(scientific)} in the natural \omterm{def:g10:scientific-biology:science}{sciences} often names a regular quantitative
pattern; in \omterm{def:g10:scientific-biology:science}{biology}, theories that organise mechanisms matter more than
equations. Neither word means ``mere guess''.
\end{definition}

\begin{proposition}[Science is self-correcting]\label{prop:g10:scientific-biology:self-correct}
Scientific knowledge is provisional: new \omterm{def:g10:scientific-biology:observation}{observations} can force revision.
A single well-designed \omterm{def:g10:scientific-biology:experiment}{experiment} that contradicts a claim is more
powerful than many that merely agree with it.
\end{proposition}
\admitted

\begin{method}[Reading a graph of experimental data]\label{met:g10:scientific-biology:graph}
\begin{enumerate}
  \item Identify axes: independent on the horizontal, dependent on the
        vertical, with units.
  \item Note scale, range and whether points are raw data or means.
  \item Describe the pattern in words (increase, plateau, threshold)
        before inventing a story.
  \item Ask what was controlled --- and what was not measured.
\end{enumerate}
\end{method}

\begin{omfigure}
\begin{tikzpicture}[font=\small]
  \begin{axis}[
    width=8.5cm, height=5.2cm,
    axis lines=left,
    xlabel={light intensity (arbitrary units)},
    ylabel={mean growth rate (\unit{mm/d})},
    xmin=0, xmax=10.5, ymin=0, ymax=5.5,
    xtick={0,2,4,6,8,10},
    ytick={0,1,2,3,4,5},
    tick label style={font=\footnotesize},
    label style={font=\small},
  ]
    \addplot[only marks, mark=*, omDef, mark size=2.2pt]
      coordinates {(0,0.4) (2,1.5) (4,2.6) (6,3.4) (8,3.9) (10,4.1)};
    \addplot[omDef, thick, smooth, tension=0.6]
      coordinates {(0,0.4) (2,1.5) (4,2.6) (6,3.4) (8,3.9) (10,4.1)};
  \end{axis}
\end{tikzpicture}

{\small \omterm{def:g10:scientific-biology:properties}{Growth} rate rises with light then flattens: more light is not
always more \omterm{def:g10:scientific-biology:properties}{growth}. Axes carry units; points are means of replicates.}
\end{omfigure}

\section{Properties of life}

What counts as alive? No single trait is enough, but a short list of
shared features covers \omterm{def:g10:scientific-biology:organism}{organisms} from bacteria to blue whales.

\begin{definition}[Organism]\label{def:g10:scientific-biology:organism}
An \emph{organism}\index{organism} is an individual living system that
maintains itself, exchanges matter and energy with its surroundings,
and can grow and reproduce (or belongs to a lineage that does).
\end{definition}

\begin{definition}[Properties of life]\label{def:g10:scientific-biology:properties}
Living \omterm{def:g10:scientific-biology:organism}{organisms} typically show:
\begin{enumerate}
  \item \emph{Organisation}\index{organisation (of life)} --- ordered
        structure from molecules to cells (and often tissues and organs);
  \item \emph{Metabolism}\index{metabolism} --- controlled chemical
        reactions that transform energy and materials;
  \item \emph{Homeostasis}\index{homeostasis} --- regulation that keeps
        internal conditions within workable limits;
  \item \emph{Growth and development}\index{growth}\index{development}
        --- increase in size and/or orderly change of form;
  \item \emph{Reproduction}\index{reproduction} --- production of new
        \omterm{def:g10:scientific-biology:organism}{organisms} of the same kind (sexually or asexually);
  \item \emph{Response}\index{response (to stimulus)} to stimuli from
        the environment;
  \item \emph{Evolutionary adaptation}\index{adaptation!evolutionary}
        --- populations change over generations under natural selection.
\end{enumerate}
Viruses sit on the edge: they have genetic material and evolve, but
cannot metabolise or reproduce without a host cell.
\end{definition}

\begin{example}[Seed versus stone]\label{ex:g10:scientific-biology:seed}
A dry seed looks inert, yet it is an \omterm{def:g10:scientific-biology:organism}{organism} in a \omterm{def:g10:scientific-biology:properties}{low-metabolism}
state: water and warmth restart respiration and \omterm{def:g10:scientific-biology:properties}{growth}. A stone may
``grow'' mineral layers, but it has no \omterm{def:g10:scientific-biology:properties}{metabolism}, no heredity and no
regulated \omterm{def:g10:scientific-biology:properties}{development}.
\end{example}

\begin{proposition}[Cell as the unit of life]\label{prop:g10:scientific-biology:cell-unit}
All known free-living \omterm{def:g10:scientific-biology:organism}{organisms} are made of one or more cells, and new
cells arise only from existing cells. This \emph{cell
theory}\index{cell theory} is developed in \cref{ch:g10:the-cell}.
\end{proposition}
\admitted

\begin{remark}[Levels of biological organisation]\label{rem:g10:scientific-biology:levels}
\omterm{def:g10:scientific-biology:science}{Biology} moves among scales: molecule $\to$ organelle $\to$ cell $\to$
tissue $\to$ organ $\to$ \omterm{def:g10:scientific-biology:organism}{organism} $\to$ population $\to$ community $\to$
ecosystem $\to$ biosphere. A good question names its scale.
\end{remark}

\begin{omfigure}
\begin{tikzpicture}[font=\footnotesize, node distance=0.15cm,
  L/.style={draw, rounded corners, fill=omDef!8, minimum width=2.1cm,
    minimum height=0.55cm, align=center, inner sep=2pt}]
  \node[L] (m) at (0,0) {molecule};
  \node[L] (o) at (2.4,0) {organelle};
  \node[L] (c) at (4.8,0) {cell};
  \node[L] (t) at (7.2,0) {tissue};
  \node[L] (og) at (0,-1.1) {organ};
  \node[L] (or) at (2.4,-1.1) {organism};
  \node[L] (p) at (4.8,-1.1) {population};
  \node[L] (cm) at (7.2,-1.1) {community};
  \node[L, fill=omMeth!12, minimum width=4.5cm] (e) at (3.6,-2.2)
    {ecosystem $\to$ biosphere};
  \foreach \a/\b in {m/o,o/c,c/t,t/og,og/or,or/p,p/cm}
    \draw[-Stealth, omDef] (\a) -- (\b);
  \draw[-Stealth, omMeth] (cm.south) -- ++(0,-0.25) -| (e.north);
\end{tikzpicture}

{\small Nested levels of organisation: each level has properties the
parts alone do not show.}
\end{omfigure}

\section{Exercises}

\begin{exercise}[$\star$]\label{exo:g10:scientific-biology:1}
For each statement, say whether it is primarily an \emph{\omterm{def:g10:scientific-biology:observation}{observation}}
or an \emph{\omterm{def:g10:scientific-biology:observation}{inference}}:
\begin{enumerate}
  \item ``The pond water contains green filaments visible under the
        microscope.''
  \item ``Those filaments are photosynthetic bacteria.''
  \item ``After \qty{48}{h} in the dark, the filaments look brown.''
  \item ``Darkness killed the photosynthetic machinery.''
\end{enumerate}
\end{exercise}

\begin{exercise}[$\star$]\label{exo:g10:scientific-biology:2}
Which of the following are testable scientific hypotheses? Rewrite any
that are not so that they become testable.
\begin{enumerate}
  \item Cats are more intelligent than dogs.
  \item Increasing temperature from \qty{15}{\celsius} to
        \qty{25}{\celsius} doubles the hopping rate of a given frog
        species under otherwise fixed conditions.
  \item Life is beautiful.
  \item Seedlings watered with \qty{1}{\percent} salt solution grow
        less tall in \qty{10}{d} than seedlings watered with pure water.
\end{enumerate}
\end{exercise}

\begin{exercise}[$\star$]\label{exo:g10:scientific-biology:3}
In an \omterm{def:g10:scientific-biology:experiment}{experiment} on caffeine and heart rate in water fleas, identify
the \omterm{def:g10:scientific-biology:variables}{independent variable}, the \omterm{def:g10:scientific-biology:variables}{dependent variable}, and three controlled
variables you would insist on.
\end{exercise}

\begin{exercise}[$\star$]\label{exo:g10:scientific-biology:4}
Why is a \omterm{def:g10:scientific-biology:control}{control group} necessary? Invent a control for a trial of a new
plant fertiliser sold as ``doubling tomato yield''.
\end{exercise}

\begin{exercise}[$\star$]\label{exo:g10:scientific-biology:5}
List five properties of life from \cref{def:g10:scientific-biology:properties}
and give a concrete example of each in a houseplant.
\end{exercise}

\begin{exercise}[$\star\star$]\label{exo:g10:scientific-biology:6}
A student claims: ``I proved that music makes plants grow because my
plant next to the speaker grew taller than the plant in the quiet
room.'' Critique the design using \cref{met:g10:scientific-biology:fair-test}.
List at least four weaknesses and how to fix each.
\end{exercise}

\begin{exercise}[$\star\star$]\label{exo:g10:scientific-biology:7}
Sketch axes (labels and units) for a graph of mean leaf area versus
daily light duration in hours for lettuce seedlings. Where would you
plot the control? What pattern would support the idea that light
duration limits \omterm{def:g10:scientific-biology:properties}{growth} up to a plateau?
\end{exercise}

\begin{exercise}[$\star\star$]\label{exo:g10:scientific-biology:8}
Viruses are sometimes called ``borderline living''. Using the list of
life properties, argue both sides in a short paragraph, then state your
own working conclusion for this course.
\end{exercise}

\begin{exercise}[$\star\star$]\label{exo:g10:scientific-biology:9}
Distinguish \emph{\omterm{def:g10:scientific-biology:hypothesis}{hypothesis}}, \emph{\omterm{def:g10:scientific-biology:theory}{theory}} and \emph{\omterm{def:g10:scientific-biology:theory}{law}} as used in
this chapter. Place \omterm{prop:g10:scientific-biology:cell-unit}{cell theory} and ``all \omterm{def:g10:scientific-biology:organism}{organisms} are made of cells''
in that framework.
\end{exercise}

\begin{exercise}[$\star\star\star$]\label{exo:g10:scientific-biology:10}
Design an \omterm{def:g10:scientific-biology:experiment}{experiment} to test whether earthworms prefer moist soil to
dry soil. Include \omterm{def:g10:scientific-biology:hypothesis}{hypothesis}, variables, control or comparison, sample
size, and what result would support or reject the \omterm{def:g10:scientific-biology:hypothesis}{hypothesis}.
\end{exercise}

\begin{exercise}[$\star\star\star$]\label{exo:g10:scientific-biology:11}
A paper reports that a drug ``significantly improves memory'' in mice,
with $n=4$ per group and no blinding. What questions would a careful
reader ask before trusting the claim?
\end{exercise}

\section{Problem: Designing a schoolyard investigation}

\begin{problem}[{Weekend problem --- from a vague question to a
protocol: plan a full investigation of lichen cover on tree bark as a
rough air-quality indicator}]
\label{pb:g10:scientific-biology:1}
Your class wants to compare lichen abundance on tree trunks near a busy
road and in a quieter park \qty{800}{m} away. Lichens are known to be
sensitive to some air pollutants, but many other factors also matter.

\medskip
\noindent\textbf{Part I --- Framing.}
\begin{enumerate}
  \item Write a focused research question.
  \item State a testable \omterm{def:g10:scientific-biology:hypothesis}{hypothesis} with a clear prediction.
  \item Name independent and \omterm{def:g10:scientific-biology:variables}{dependent variables} (with units or scoring
        method). List at least five controlled or matched variables.
\end{enumerate}

\medskip
\noindent\textbf{Part II --- Protocol.}
\begin{enumerate}
  \item Describe how you will choose trees and sample bark (height,
        aspect, number of trees, replicates).
  \item Explain why a single tree per site would be inadequate.
  \item Identify two ethical or safety constraints of fieldwork.
\end{enumerate}

\medskip
\noindent\textbf{Part III --- Interpretation.}
\begin{enumerate}
  \item Suppose mean lichen cover is higher in the park. List three
        alternative explanations that are \emph{not} air quality.
  \item What additional measurements would help separate those
        explanations?
  \item In one paragraph, explain why this study is observational
        rather than a controlled \omterm{def:g10:scientific-biology:experiment}{experiment}, and what that limits.
\end{enumerate}
\end{problem}
